\chapter{Expectation}
\label{ch:expectation}

\begin{question}
What is the average value of a random variable, how do we compute it from its distribution, and what does it tell us about the probability of large values?
\end{question}

\begin{whatfor}
\begin{itemize}
  \item Defines the expectation $\E_P(\xi)=\int_\Omega\xi\,dP$, built from indicators and step functions.
  \item Gives the working formula $\E_P h(\xi)=\int h(x)f_\xi(x)\,dx$ (or a sum), then variance and covariance.
  \item Computes the normal moments and the normal exponential moment $\E_P e^{\xi}=e^{\mu+\sigma^2/2}$, used all through the book.
  \item Proves the Markov and Chebyshev inequalities.
\end{itemize}
\skipline{fs:expectation}
\end{whatfor}

\begin{prereq}
From \crefrange{ch:spaces}{ch:rv}: probability space, random variable, indicator, distribution, density, discrete distribution, the normal distribution and $\Phi$.
\end{prereq}

%=====================================================================
\section{The expectation as an integral}

Throughout, $(\Omega,\F,P)$ is a probability space and $\E_P$ is the average under $P$.

\begin{definition}[Integrable; expectation]
\label{def:expectation}
A random variable $\xi$ is \dfn{integrable} if $\int_\Omega|\xi|\,dP<\infty$. Then its \dfn{expectation} is
\[
  \E_P(\xi)=\int_\Omega\xi\,dP,
\]
the average of the values $\xi(\omega)$, each weighted by its probability. The integrable random variables form the family $L^1=L^1(\Omega,\F,P)$.
\end{definition}

The integral is built in three steps, each one extending the last.

\begin{steps}
  \item \emph{Indicators.} For an event $A$, $\E_P(\ind{A})=P(A)$: the value is $1$ on $A$ and $0$ elsewhere, so only the probability of $A$ survives.
  \item \emph{Finitely many values.} A \dfn{step function} is $\eta=\sum_{i=1}^n x_i\,\ind{A_i}$ with numbers $x_i$ and pairwise disjoint events $A_i$; it takes the value $x_i$ on $A_i$. Its expectation is
  \[
    \E_P(\eta)=\sum_{i=1}^n x_i\,P(A_i).
  \]
  \item \emph{Everything else.} Any $\xi\ge0$ is the limit of step functions $0\le\eta_1\le\eta_2\le\dots\uparrow\xi$, and $\E_P(\xi)=\lim_n\E_P(\eta_n)$ (the \dfn{monotone convergence theorem}). A general $\xi$ is split as $\xi=\xi^+-\xi^-$ with $\xi^+=\max(\xi,0)$, $\xi^-=\max(-\xi,0)$, and $\E_P(\xi)=\E_P(\xi^+)-\E_P(\xi^-)$.
\end{steps}

\begin{example}[Rounding a uniform point down]
\label{ex:dyadic}
$U$ uniform on $[0,1]$. Round it down to a multiple of $2^{-n}$: $\eta_n=\lfloor2^nU\rfloor/2^n$. Then $\eta_n$ takes the $2^n$ values $k/2^n$, each with probability $2^{-n}$, so
\[
  \E_P(\eta_n)=\frac{1}{2^n}\sum_{k=0}^{2^n-1}\frac{k}{2^n}=\frac{2^n-1}{2^{n+1}}=\tfrac12\big(1-2^{-n}\big)\;\uparrow\;\tfrac12=\E_P(U).
\]
For $n=1,2,3$: $0.25$, $0.375$, $0.4375$.
\end{example}

\begin{pic}
\begin{tikzpicture}
\begin{axis}[width=8.5cm,height=4.6cm,xmin=0,xmax=1,ymin=0,ymax=1.05,xlabel={$\omega$ (the value of $U$)},ylabel={value},xtick={0,0.25,0.5,0.75,1},ytick={0,0.25,0.5,0.75,1},legend style={font=\scriptsize,at={(0.98,0.02)},anchor=south east}]
  \addplot[black,thick,domain=0:1] {x}; \addlegendentry{$U$}
  \addplot[defc,very thick,const plot] coordinates {(0,0)(0.5,0.5)(1,0.5)}; \addlegendentry{$\eta_1$}
  \addplot[thmc,very thick,const plot] coordinates {(0,0)(0.25,0.25)(0.5,0.5)(0.75,0.75)(1,0.75)}; \addlegendentry{$\eta_2$}
\end{axis}
\end{tikzpicture}
\picnote{Step functions climb up to $U$ from below. Their expectations are simple sums, and they increase to $\E_P(U)=\frac12$.}
\end{pic}

The absolute value in the definition matters. If $\E_P(\xi^+)=\E_P(\xi^-)=\infty$, then ``$\infty-\infty$'' has no value. Example: $P(\xi=2^k)=2^{-k}$ for $k=1,2,\ldots$ gives $\E_P|\xi|=\sum 2^k2^{-k}=1+1+\cdots=\infty$, so $\xi$ is not integrable.

\begin{result}[Rules for expectation]
\label{res:E-rules}
For integrable $\xi,\eta$ and numbers $a,b$:
\begin{enumerate}[label=(\alph*)]
  \item linearity: $\E_P(a\xi+b\eta)=a\E_P(\xi)+b\E_P(\eta)$;
  \item monotonicity: $\xi\le\eta$ $\Rightarrow$ $\E_P(\xi)\le\E_P(\eta)$;
  \item constants: $\E_P(c)=c$.
\end{enumerate}
\end{result}
\begin{proof}
For step functions these are rules for finite sums: put both on common disjoint events and add. They pass to limits by the construction.
\end{proof}

%=====================================================================
\section{Computing an expectation from the distribution}

We never integrate over $\Omega$ in practice. We move to the real line, where the distribution lives.

\begin{theorem}[Expectation of a function of $\xi$]
\label{thm:lotus}
Let $h$ be a Borel function with $h(\xi)$ integrable. Then $\E_P\,h(\xi)=\int_\R h(x)\,dP_\xi(x)$. In the two working cases:
\[
  \E_P\,h(\xi)=\int_{-\infty}^{\infty}h(x)\,f_\xi(x)\,dx\quad\text{(density)},\qquad
  \E_P\,h(\xi)=\sum_i h(x_i)\,P(\xi=x_i)\quad\text{(discrete)}.
\]
\end{theorem}
\begin{proof}
Prove it for step functions $h$; the rest follows by the same limits as the construction.
\begin{steps}
  \item Let $h=\sum_i h_i\ind{A_i}$ with disjoint Borel sets $A_i$. Then $h(\xi)=\sum_ih_i\,\ind{\xi\in A_i}$, a step function on $\Omega$.
  \item So $\E_P\,h(\xi)=\sum_ih_i\,P(\xi\in A_i)=\sum_ih_i\,P_\xi(A_i)=\int_\R h\,dP_\xi$. \qedhere
\end{steps}
\end{proof}

\begin{example}[Two ways, one answer]
Two fair tosses, $\xi$ = number of heads, $h(x)=x^2$.
\begin{itemize}
  \item Over outcomes: $HH\to4$, $HT\to1$, $TH\to1$, $TT\to0$, so $\E_P(\xi^2)=\frac14(4+1+1+0)=1.5$.
  \item Over values: $0^2\cdot\frac14+1^2\cdot\frac12+2^2\cdot\frac14=1.5$.
\end{itemize}
\end{example}

In general $\E_P\,h(\xi)\ne h(\E_P\xi)$: apply $h$ to each value, then average. Die: $\E_P(1/\xi)=\frac16(1+\frac12+\dots+\frac16)=0.408$, while $1/\E_P(\xi)=1/3.5=0.286$.

\begin{definition}[Mean, variance, covariance]
\label{def:var}
For $\xi$ with $\E_P(\xi^2)<\infty$, the \textbf{mean}\index{mean} is $\E_P(\xi)$ and the \dfn{variance} is
\[
  \Var(\xi)=\E_P\big[(\xi-\E_P\xi)^2\big]=\E_P(\xi^2)-(\E_P\xi)^2 .
\]
For two such random variables, the \dfn{covariance} is $\Cov(\xi,\eta)=\E_P\big[(\xi-\E_P\xi)(\eta-\E_P\eta)\big]=\E_P(\xi\eta)-\E_P(\xi)\E_P(\eta)$.
\end{definition}
The second forms follow by expanding the square and using linearity. Variance and covariance are always taken under $P$ in Parts~0--3, so we write them without a subscript.

\begin{result}[Variance and covariance rules]
\label{res:var-rules}
For numbers $a,b,c$ and random variables with finite second moments:
\begin{enumerate}[label=(\alph*)]
  \item $\Var(a\xi+c)=a^2\Var(\xi)$;
  \item $\Var\big(\sum_i a_i\xi_i\big)=\sum_i\sum_j a_ia_j\Cov(\xi_i,\xi_j)$; for two: $\Var(a\xi+b\eta)=a^2\Var\xi+b^2\Var\eta+2ab\Cov(\xi,\eta)$;
  \item $\Cov$ is linear in each slot, $\Cov(\xi,\xi)=\Var(\xi)$, and $\Cov(\xi,c)=0$;
  \item $\big(\E_P|\xi\eta|\big)^2\le\E_P(\xi^2)\,\E_P(\eta^2)$, and so $\big(\E_P|\xi|\big)^2\le\E_P(\xi^2)$ \;(\dfn{Cauchy--Schwarz}).
\end{enumerate}
\end{result}
\begin{proof}
(a)--(c) expand the products and use linearity. For (d): for every number $t$, $0\le\E_P\big[(|\xi|-t|\eta|)^2\big]=\E_P\xi^2-2t\,\E_P|\xi\eta|+t^2\E_P\eta^2$. A quadratic in $t$ that is never negative has discriminant $\le0$: $4(\E_P|\xi\eta|)^2-4\E_P\xi^2\,\E_P\eta^2\le0$. Take $\eta=1$ for the last claim.
\end{proof}

%=====================================================================
\section{Moments of the normal distribution}

\begin{result}[Normal moments]
\label{res:normal-moments}
Let $Z\sim N(0,1)$ and $\xi=\mu+\sigma Z\sim N(\mu,\sigma^2)$.
\begin{enumerate}[label=(\alph*)]
  \item $\E_P(Z)=0$, $\E_P(Z^2)=1$, $\E_P(Z^3)=0$, $\E_P(Z^4)=3$. So $\E_P(\xi)=\mu$ and $\Var(\xi)=\sigma^2$.
  \item \dfn{Exponential moment}: for every number $u$,
  \[
    \boxed{\;\E_P\big(e^{uZ}\big)=e^{u^2/2},\qquad \E_P\big(e^{u\xi}\big)=e^{u\mu+\frac12u^2\sigma^2}.\;}
  \]
  In particular $\E_P(e^{\xi})=e^{\mu+\sigma^2/2}$.
\end{enumerate}
\end{result}
\begin{proof}
Write $\varphi(x)=\frac{1}{\sqrt{2\pi}}e^{-x^2/2}$ for the standard normal density.
\begin{steps}
  \item \emph{Odd moments.} $x\varphi(x)$ and $x^3\varphi(x)$ are odd functions, so their integrals over $\R$ are $0$.
  \item \emph{Even moments.} Since $\varphi'(x)=-x\varphi(x)$, integration by parts gives $\int x^{k}\varphi\,dx=\int x^{k-1}\cdot x\varphi\,dx=(k-1)\int x^{k-2}\varphi\,dx$. So $\E_P Z^2=1\cdot1=1$ and $\E_P Z^4=3\cdot\E_P Z^2=3$.
  \item \emph{Complete the square.} $ux-\frac{x^2}{2}=\frac{u^2}{2}-\frac{(x-u)^2}{2}$, so
  \[
    \E_P\big(e^{uZ}\big)=\int_{-\infty}^{\infty}\frac{1}{\sqrt{2\pi}}e^{ux-x^2/2}\,dx
    =e^{u^2/2}\int_{-\infty}^{\infty}\frac{1}{\sqrt{2\pi}}e^{-(x-u)^2/2}\,dx=e^{u^2/2},
  \]
  because the last integral is the total mass of the $N(u,1)$ density.
  \item \emph{General normal.} $\E_P(e^{u\xi})=e^{u\mu}\,\E_P\big(e^{(u\sigma)Z}\big)=e^{u\mu}e^{u^2\sigma^2/2}$. Linearity gives $\E_P\xi=\mu$; and $\Var(\xi)=\sigma^2\Var(Z)=\sigma^2$. \qedhere
\end{steps}
\end{proof}

\begin{example}[The mean of $e^{\xi}$]
\label{ex:lognormal}
$\xi\sim N(0.1,\,0.04)$, so $\mu=0.1$, $\sigma^2=0.04$. Then $\E_P(e^{\xi})=e^{0.1+0.02}=e^{0.12}\approx1.1275$, larger than $e^{\E_P\xi}=e^{0.1}\approx1.1052$. The extra $\frac12\sigma^2$ is the price of randomness inside a convex function; it returns in the stock-price models of Parts~2 and~4.
\end{example}

%=====================================================================
\section{Markov and Chebyshev inequalities}

\begin{theorem}[Markov and Chebyshev]
\label{thm:markov}
Two bounds on tail probabilities.
\begin{enumerate}[label=(\alph*)]
  \item \dfn{Markov inequality}: if $X\ge0$ and $x>0$, then $P(X\ge x)\le\dfrac{\E_P(X)}{x}$.
  \item \dfn{Chebyshev inequality}: if $\xi$ has mean $\mu$ and variance $\sigma^2$, then for every $\varepsilon>0$, $P(|\xi-\mu|\ge\varepsilon)\le\dfrac{\sigma^2}{\varepsilon^2}$.
\end{enumerate}
\end{theorem}
\begin{proof}
One picture does both.
\begin{steps}
  \item $X\ge x\,\ind{X\ge x}$: where $X\ge x$ the right side is $x\le X$; elsewhere it is $0\le X$.
  \item Take expectations (monotonicity): $\E_P(X)\ge x\,P(X\ge x)$. Divide by $x$. This is (a).
  \item For (b), apply (a) to $X=(\xi-\mu)^2\ge0$ with $x=\varepsilon^2$: $P\big((\xi-\mu)^2\ge\varepsilon^2\big)\le\frac{\E_P(\xi-\mu)^2}{\varepsilon^2}=\frac{\sigma^2}{\varepsilon^2}$, and $\{(\xi-\mu)^2\ge\varepsilon^2\}=\{|\xi-\mu|\ge\varepsilon\}$. \qedhere
\end{steps}
\end{proof}

\begin{example}[How loose is Markov?]
\label{ex:markov}
Two comparisons of the exact tail with the Markov bound.
\begin{itemize}
  \item $X$ exponential with mean $1$: $P(X>x)=e^{-x}$ against $\E_P(X)/x=1/x$. At $x=\frac12$: $e^{-1/2}\approx0.607$ against $2$ (useless). At $x=3$: $0.050$ against $0.333$.
  \item $X$ Bernoulli$(p)$, $x=1$: $P(X\ge1)=p$ and $\E_P(X)/1=p$. Here Markov is exact.
\end{itemize}
\end{example}

\begin{pic}
\begin{tikzpicture}
\begin{axis}[width=8.5cm,height=4.4cm,xmin=0.3,xmax=5,ymin=0,ymax=1.6,xlabel={$x$},legend style={font=\scriptsize},samples=100]
  \addplot[defc,very thick,domain=0.3:5] {exp(-x)}; \addlegendentry{$P(X>x)=e^{-x}$}
  \addplot[trapc,very thick,dashed,domain=0.62:5] {1/x}; \addlegendentry{Markov bound $1/x$}
\end{axis}
\end{tikzpicture}
\picnote{The Markov bound always lies above the true tail, but it can be far above it: it only uses the mean.}
\end{pic}

%=====================================================================
\formulasheet{fs:expectation}

\begin{mastercard}[label=card:Eh]{Expectation of any function of a random variable}
Apply $h$ to each value, weight by its probability (or density), and add (or integrate). It exists when $\E_P|h(\xi)|<\infty$:
\[
  \boxed{\begin{gathered}\E_P\,h(\xi)=\sum_ih(x_i)\,P(\xi=x_i)\ \ \text{or}\ \ \int_{-\infty}^{\infty}h(x)f_\xi(x)\,dx;\\ \E_P\,h(\xi,\eta)=\iint h(x,y)f(x,y)\,dx\,dy.\end{gathered}}
\]
\begin{plugrows}
  \plugrow{Probability as an expectation}{$h=\ind{x\in A}$}{$\E_P\ind{\xi\in A}=P(\xi\in A)$}{always}{die: $\E_P\ind{\xi\ge5}=\frac13$}
  \plugrow{A payoff with finitely many values}{$h$ = value $x_i$ on event $A_i$ (disjoint)}{$\sum_ix_iP(A_i)$}{always}{$2\ind{U<0.3}-\ind{U\ge0.7}$: $0.6-0.3=0.3$}
  \plugrow{Mean, second moment, variance}{$h=x$; $h=x^2$; $h=(x-\E_P\xi)^2$}{$\E_P\xi$; $\E_P\xi^2$; $\Var\xi=\E_P\xi^2-(\E_P\xi)^2$}{integrable}{$U$ on $[0,1]$: $\frac12$, $\frac13$, $\frac1{12}$}
  \plugrow{Payoff above a level $K$}{$h=(x-K)^+$}{$\int_K^\infty(x-K)f(x)dx$}{integrable}{Exp$(1)$, $K=1$: $\int_1^\infty(x-1)e^{-x}dx=e^{-1}$}
  \plugrow{Is $\xi$ integrable?}{$h=|x|$}{finite or not}{needed before $\E_P\xi$ exists}{$P(\xi=2^k)=2^{-k}$: $\sum2^k2^{-k}=\infty$, not integrable}
  \plugrow{A function is not averaged by its value}{compare $\E_Ph(\xi)$ with $h(\E_P\xi)$}{different in general}{only equal for linear $h$}{die: $\E_P\frac1\xi=0.408$ but $\frac1{3.5}=0.286$}
\end{plugrows}
\end{mastercard}

\begin{mastercard}[label=card:var-lin]{Variance of any linear combination}
Expand the square: every pair contributes its covariance, with $\Cov(X,X)=\Var X$:
\[
  \boxed{\begin{gathered}\Var\Big(\sum_{i=1}^na_iX_i+c\Big)=\sum_{i=1}^n\sum_{j=1}^na_ia_j\Cov(X_i,X_j),\\ \Cov(X,Y)=\E_P(XY)-\E_PX\,\E_PY.\end{gathered}}
\]
\begin{plugrows}
  \plugrow{One variable}{$n=1$}{$\Var(aX+c)=a^2\Var X$ (the constant drops)}{always}{$\Var X=2$: $\Var(3X+7)=18$}
  \plugrow{Two variables}{$n=2$}{$a^2\Var X+b^2\Var Y+2ab\Cov(X,Y)$}{always}{$\Var X=2$, $\Var Y=3$, $\Cov=1$: $\Var(2X-Y)=8+3-4=7$}
  \plugrow{Covariance of combinations}{bilinearity}{$\Cov(\sum a_iX_i,\sum b_jY_j)=\sum_{i,j}a_ib_j\Cov(X_i,Y_j)$}{always}{$\Cov(X+Y,X-Y)=\Var X-\Var Y$}
  \plugrow{Uncorrelated terms}{all $\Cov(X_i,X_j)=0$, $i\ne j$}{$\sum a_i^2\Var X_i$}{zero covariances; Chapter~4 gives the common case}{sum of $n$ fair dice: $n\cdot\frac{35}{12}$}
  \plugrow{A bound (Cauchy--Schwarz)}{any two variables}{$(\E_P|XY|)^2\le\E_PX^2\,\E_PY^2$; so $(\E_P|X|)^2\le\E_PX^2$}{finite second moments}{$X\sim N(0,\sigma^2)$: $\E_P|X|\le\sigma$ (exactly $0.80\sigma$)}
\end{plugrows}
\end{mastercard}

\begin{mastercard}[label=card:normal-mgf]{Exponential moment of a normal (all normal moments)}
For $\xi\sim N(\mu,\sigma^2)$ and any number $u$:
\[
  \boxed{\begin{gathered}\E_P\,e^{u\xi}=\exp\Big(u\mu+\tfrac12u^2\sigma^2\Big).\end{gathered}}
\]
Expanding $\E_Pe^{u(\xi-\mu)}=e^{u^2\sigma^2/2}$ in powers of $u$ gives every central moment.
\begin{plugrows}
  \plugrow{$\E_Pe^{\xi}$, $\E_Pe^{a\xi+b}$}{$u=1$; $u=a$ and multiply by $e^b$}{$e^{\mu+\sigma^2/2}$; $e^{b+a\mu+a^2\sigma^2/2}$}{always}{$N(0.1,0.04)$: $\E_Pe^\xi=e^{0.12}\approx1.1275$}
  \plugrow{Variance of $e^\xi$}{$u=2$ minus the square of $u=1$}{$e^{2\mu+\sigma^2}(e^{\sigma^2}-1)$}{always}{$N(0,1)$: $e(e-1)\approx4.67$}
  \plugrow{Central moments}{match powers of $u$ in $e^{u^2\sigma^2/2}=1+\frac{u^2\sigma^2}{2}+\frac{u^4\sigma^4}{8}+\cdots$}{$\E_P(\xi-\mu)^2=\sigma^2$, $\E_P(\xi-\mu)^4=3\sigma^4$, odd ones $0$}{the coefficient of $u^k$ equals $\E_P(\xi-\mu)^k/k!$}{$N(0,2)$: $\E_P\xi^4=3\cdot4=12$}
  \plugrow{Raw moments}{expand $\xi=\mu+(\xi-\mu)$}{$\E_P\xi=\mu$, $\E_P\xi^2=\mu^2+\sigma^2$}{always}{$N(1,4)$: $\E_P\xi^2=5$}
\end{plugrows}
\end{mastercard}

\begin{mastercard}[label=card:markov]{Bounding a tail by an expectation}
For any function $g\ge0$ and number $c>0$:
\[
  \boxed{\begin{gathered}P\big(g(X)\ge c\big)\le\frac{\E_P\,g(X)}{c}.\end{gathered}}
\]
\begin{plugrows}
  \plugrow{Markov}{$g(x)=x$}{$P(X\ge c)\le\E_PX/c$}{$X\ge0$}{Exp, mean 1, $c=3$: exact $0.050$, bound $0.333$}
  \plugrow{Chebyshev}{$g(x)=(x-\mu)^2$, $c=\varepsilon^2$}{$P(|X-\mu|\ge\varepsilon)\le\sigma^2/\varepsilon^2$}{finite variance}{$\sigma=1$, $\varepsilon=2$: $\le\frac14$}
  \plugrow{When is it exact?}{$X$ takes only the values $0$ and $c$}{equality}{$g(X)$ is $0$ or exactly $c$}{Bernoulli$(p)$, $c=1$: both sides $p$}
\end{plugrows}
\end{mastercard}
